% =============================================================================
% §2 Related Work (~0.7 page) — narrative, minimal redundancy with §1
% =============================================================================
\section{Related Work}
\label{sec:related}

\subsection{Skill acquisition and curation}

Work on the skill layer divides by when skills enter
the system.  One strand acquires them at run time, while the
agent is solving tasks.  Voyager grows a skill library through
environment exploration \citep{Wang2023Voyager}, Toolformer
trains the model to invoke tools self-supervisedly
\citep{Schick2023Toolformer}, ToolLLM and Gorilla scale API
calling to large tool inventories
\citep{Qin2024ToolLLM,Patil2023Gorilla}, and recent
agent-evolution work lets agents redesign their own scaffolding
\citep{Zhang2024Memento,AEvolve2025}.  All of these build on
agent harness scaffolds \citep{Yao2023ReAct,Shinn2023Reflexion}.
The other strand, which this paper occupies, prepares skills
before the agent ever runs.  We curate a
community-contributed corpus at ingest time, so that runtime
systems have a trustworthy library to draw from.  This
ingest-time curation mirrors data-quality work in language-model
training, where deduplication \citep{Lee2022Dedup} and
quality-based data selection \citep{Zhou2023LIMA} improve
downstream models more than raw scale.

Released community skill corpora vary in emphasis.
\citet{SkillNet2026}, concurrent with our work, releases a
200K-skill ontology rated along five quality rubrics and linked
by a skill-relation graph.  It measures benefit on
simulated environments, with skills synthesised from each
benchmark's own expert trajectories rather than on real tasks
served by a deployed corpus.  \citet{SkillFlow2026} pairs a 36K
uncurated library with a four-stage retrieval pipeline and
treats quality only implicitly, through artifact-density
signals.
\citet{AgentSkillsDataAnalysis2026} provides a data-driven
characterisation of a smaller curated set.  None of these releases
reports a source-licence audit, so their redistribution rights are
left unresolved.  Nor does prior work tie these per-skill scores to
their downstream operational role.  We introduce a
compact three-facet framework chosen for measurable facet
separation rather than rubric coverage, and characterise its
operational signal directly (\S\ref{sec:claim-a}).

\subsection{Skill retrieval and evaluation}

Skill retrieval architectures range from lexical and dense
recall to learned skill routers \citep{SkillRouter2026},
supported by dedicated retrieval benchmarks \citep{SkillRet2026},
corpus-size scaling analyses \citep{ScalingLaws2026}, and
dependency-aware composition \citep{GraphOfSkills2026}.  A
complementary line of work emphasises skill orchestration and
context-injection over retrieval per se.
\citet{AgentSkillOS2026} couples capability-tree retrieval with
DAG-based orchestration at $\sim$$200$K scale, and
\citet{SkillRAE2026} reports $+11.7\%$ on SkillsBench via a
skill graph plus compact compilation.  That gain, however, is
measured against the prior state-of-the-art injection method
rather than a no-skill or community-corpus baseline, so it is not
directly comparable to our end-to-end $\Delta$.
The retrieval strand reports ranking metrics on its
own curated pools rather than task outcomes.  Our
stack is fine-tuned on the deduplicated
corpus it serves and adds an LLM selector that reads full skill
bodies before injection, with both its standalone retrieval
performance and its end-to-end effect on task pass rates
measured directly.

On evaluation, \citet{SkillsBench2026} characterises oracle-skill
efficacy, and \citet{UCSBSkillsWild2026} reports diminishing
returns on uncurated libraries at scale.  Neither
evaluates a single community-curated corpus under a
deployable retrieval-and-selection stack across multiple
real-world benchmarks.  \S\ref{sec:experiments} fills this gap
with a cross-benchmark, cross-harness, multi-backbone measurement
of end-to-end task gains.
